\documentclass[11pt]{article}
\usepackage[margin=1in]{geometry}
\usepackage{amsmath,amssymb,amsthm,mathtools}
\usepackage{hyperref}
\usepackage{enumitem}
\usepackage{microtype}
\usepackage{mathrsfs}

\newtheorem{theorem}{Theorem}[section]
\newtheorem{lemma}[theorem]{Lemma}
\newtheorem{proposition}[theorem]{Proposition}
\newtheorem{corollary}[theorem]{Corollary}
\newtheorem{assumption}[theorem]{Assumption}
\theoremstyle{definition}
\newtheorem{definition}[theorem]{Definition}
\newtheorem{remark}[theorem]{Remark}
\newtheorem{example}[theorem]{Example}

\newcommand{\E}{\mathbb{E}}
\newcommand{\Var}{\mathrm{Var}}
\newcommand{\Cov}{\mathrm{Cov}}
\newcommand{\tr}{\mathrm{tr}}
\newcommand{\1}{\mathbf{1}}
\newcommand{\R}{\mathbb{R}}
\newcommand{\Z}{\mathbb{Z}}
\newcommand{\C}{\mathbb{C}}
\newcommand{\N}{\mathbb{N}}
\newcommand{\cH}{\mathcal{H}}
\newcommand{\cA}{\mathcal{A}}
\newcommand{\cF}{\mathcal{F}}
\newcommand{\dist}{\operatorname{dist}}
\newcommand{\supp}{\operatorname{supp}}
\newcommand{\Ran}{\operatorname{Ran}}
\newcommand{\Lip}{\operatorname{Lip}}
\newcommand{\ip}[2]{\langle #1,#2\rangle}
\newcommand{\norm}[1]{\lVert #1\rVert}

\title{Block Transfer Operators, Schur-Complement Contraction, and Spectral Gaps:\\
An Operator-Factorization Approach to Lattice Systems}
\author{Ramiro Fontes\\Quijotic Research\\\texttt{ramirofontes@gmail.com}}
\date{February 2026}

\begin{document}
\maketitle

\begin{abstract}
We develop a self-contained framework for extracting exponential clustering and spectral gaps
from strict gradient contraction of block transfer operators on lattice systems with compact
state spaces.  The mechanism rests on an operator-factorization philosophy:
decompose a lattice potential into interior and boundary degrees of freedom,
compute the Schur complement of the interior block to obtain effective boundary curvature,
and show that this curvature drives a strict $H^1$ contraction of the block transfer operator.
The key combinatorial input is a topological cancellation ($A^\top v = 0$) ensuring that the
coarse boundary mode is \emph{unscreenable}---it lies in the kernel of the adjoint
of the interior incidence matrix and therefore retains full curvature after integrating out
interior variables.

From this single local contraction estimate we derive, by deterministic propagation arguments:
(i) exponential decay of boundary influence across block layers;
(ii) exponential clustering of truncated correlations for local observables;
(iii) a positive spectral gap for the transfer-matrix Hamiltonian under a
reflection-positivity hypothesis.
We further show that when the contraction constant is uniform along a
scaling family parametrized by a lattice spacing $a\downarrow 0$, the spectral gap persists
in subsequential continuum limits obtained via the Osterwalder--Schrader reconstruction theorem.
The framework is formulated abstractly but is motivated by, and directly applicable to,
lattice gauge theories with compact gauge groups.
\end{abstract}

\tableofcontents

%======================================================================
\section{Introduction}\label{sec:intro}
%======================================================================

\subsection{Motivation and context}

A recurring theme in constructive statistical mechanics and quantum field theory is the
extraction of long-distance information (correlation decay, spectral gaps, phase structure)
from short-distance analytic control.  The classical Dobrushin--Shlosman program
\cite{DobrushinShlosman} achieves this for spin systems via contraction of conditional
expectations; the Balaban program \cite{Balaban1,Balaban2,Balaban3} achieves it for gauge
theories via iterated renormalization-group (RG) transformations.

In this paper we isolate a mechanism that sits between these traditions: a single-block
\emph{gradient contraction} estimate for a transfer operator obtained by integrating out
interior degrees of freedom, where the contraction is certified by an explicit
\emph{Schur-complement} computation.  The approach is operator-theoretic in spirit:
\begin{enumerate}[label=(\roman*)]
\item the lattice potential is decomposed into interior ($y$) and boundary ($x$) variables;
\item the Hessian factors into blocks $V_{xx}$, $V_{xy}$, $V_{yy}$;
\item a topological cancellation ($A^\top v=0$) forces $V_{xy}^{(0)}=0$ at leading order,
making the Schur complement \emph{exact};
\item the resulting boundary curvature drives heat-kernel smoothing with a strict contraction
factor $q_0<1$;
\item perturbative corrections from nonlinear (e.g.\ BCH) remainders are controlled by
a quantitative stability lemma for Schur complements.
\end{enumerate}

\subsection{Relationship to operator-factorization calculus}

This work extends the operator-factorization program developed in the companion papers
\cite{Fontes2026I,Fontes2026II,Fontes2026III,Fontes2026IV}.  Those works establish the
operator-theoretic identity
\[
(\mathrm{Id}-\E) \;=\; \delta\,\Pi\,D
\]
as a factorization of stochastic differentiation through projection ($\Pi$) and divergence
($\delta$), yielding Clark--Ocone representations, It\^o formulas, and Bakry--\'Emery
curvature estimates in increasing generality (Gaussian, rough fractional, non-semimartingale,
martingale-innovation settings).

The present paper adapts the same structural philosophy to a lattice-statistical-mechanics
setting:
\begin{center}
\begin{tabular}{lcl}
\emph{Stochastic calculus setting} && \emph{Lattice block setting}\\[3pt]
\hline\\[-6pt]
Conditional expectation $\E[\,\cdot\mid\cF]$ &$\longleftrightarrow$& Block transfer operator $T$\\[2pt]
Innovation / martingale increment &$\longleftrightarrow$& Boundary gradient $\nabla_\partial$\\[2pt]
Projection $\Pi$ onto range &$\longleftrightarrow$& Schur complement / $I-P_A$\\[2pt]
Contraction $\norm{(\mathrm{Id}-\E)f}\le q\norm{f}$ &$\longleftrightarrow$& $\norm{\nabla(Tf)}\le q\norm{\nabla f}$, $q<1$\\[2pt]
Clark--Ocone representation &$\longleftrightarrow$& Exponential clustering / mass gap
\end{tabular}
\end{center}
The unifying idea is: \emph{factor the operator, identify the contractive piece, and propagate
contraction to global conclusions}.  In the stochastic setting, the factorization yields
martingale representations and functional inequalities.  In the lattice setting, it yields
correlation decay and spectral gaps.

\subsection{Structure of the paper}

Section~\ref{sec:schur} develops the local block analysis: the topological cancellation
$A^\top v=0$, the exact and robust Schur complements, and the strict $H^1$ contraction
theorem.
Section~\ref{sec:clustering} propagates local contraction to exponential clustering
of lattice correlations via influence-decay arguments.
Section~\ref{sec:gap} extracts a positive spectral gap from clustering under a
reflection-positivity hypothesis, including persistence in continuum limits.
Section~\ref{sec:uniqueness} shows how contraction, combined with a functional equation
(such as a lattice Schwinger--Dyson identity) and boundary data, yields uniqueness of
solutions---the abstract ``$Ax=0$ implies $x=0$'' conclusion.
Section~\ref{sec:application} specializes the framework to SU(2) lattice gauge theory.

%======================================================================
\section{Block Schur Complement and Local Contraction}\label{sec:schur}
%======================================================================

\subsection{Setup: tiles, interior variables, and boundary variables}

We work on a $d$-dimensional hypercubic lattice $\Z^d$ ($d\ge 2$) with a block decomposition
into hypercubes of side length $L\ge 2$.  Fix a single block $B$.  After a suitable
gauge-fixing or spanning-tree reduction (in the gauge-theory application) or simply by the
geometry of the model, the degrees of freedom in $B$ decompose into:
\begin{itemize}[nosep]
\item \emph{boundary variables} $x\in M_\partial$, a compact Riemannian manifold (e.g.\ a product of copies of a compact Lie group),
\item \emph{interior variables} $y\in M_{\mathrm{int}}$, similarly compact.
\end{itemize}

In a small-field (``corridor'') regime, both $x$ and $y$ are coordinatized by exponential
charts, and we work in $\R^{n_\partial}\times\R^{n_i}$ with $\norm{(x,y)}\le\delta$ for a
corridor radius $\delta>0$.

The block potential has the form
\begin{equation}\label{eq:Vfull}
V(x,y) \;=\; V^{(0)}(x,y) + \Delta(x,y),
\end{equation}
where $V^{(0)}$ is a quadratic core and $\Delta$ collects higher-order and geometric
(Jacobian, BCH) corrections.

\subsection{The linearized plaquette map and the quadratic core}

Let $\omega(x,y)\in\R^{n_p}$ denote the vector of local constraint violations (plaquette
logarithms in the gauge setting, or more generally the local interaction terms) in the block.
In the corridor,
\begin{equation}\label{eq:omega}
\omega(x,y) \;=\; Ay + vx + e(x,y),
\end{equation}
where $A\in\R^{n_p\times n_i}$ is the interior incidence matrix, $v\in\R^{n_p}$ is the
boundary coupling vector, and $e$ is the nonlinear remainder.  The quadratic core is
\begin{equation}\label{eq:V0}
V^{(0)}(x,y) \;:=\; \frac{\beta c_2}{2}\,\norm{Ay+vx}^2,
\end{equation}
where $\beta>0$ is an inverse-temperature or inverse-coupling parameter and $c_2>0$ is a
model-dependent coefficient.

\subsection{Topological cancellation: $A^\top v=0$}

The following is the key combinatorial fact.

\begin{lemma}[Topological cancellation]\label{lem:Atv0}
For the standard block decomposition (e.g.\ a $2\times 2$ tile in each coordinate plane after
spanning-tree gauge fixing in the gauge setting, or any block where each interior variable
participates in exactly two local constraints with opposite orientations), we have
\[
A^\top v \;=\; 0.
\]
\end{lemma}

\begin{proof}
Each interior variable $y_k$ appears in exactly two constraint equations, with opposite signs:
there exist exactly two indices $p\ne p'$ with $A_{p,k}=+1$, $A_{p',k}=-1$, and
$A_{q,k}=0$ for $q\notin\{p,p'\}$.  If $v=(1,\ldots,1)^\top$ (uniform boundary coupling), then
\[
(A^\top v)_k \;=\; \sum_{p=1}^{n_p} A_{p,k}\,v_p \;=\; (+1)+(-1) \;=\; 0
\]
for every $k$.  The result holds more generally whenever $v$ lies in the left kernel of the
signed incidence structure.
\end{proof}

\begin{corollary}[Unscreenable boundary mode]\label{cor:unscreen}
Let $P_A := A(A^\top A)^{-1}A^\top$ be the orthogonal projector onto $\Ran(A)$.  Then
$P_A v = 0$, so
\[
(I - P_A)\,v \;=\; v,\qquad v^\top(I-P_A)\,v \;=\; \norm{v}^2.
\]
\end{corollary}

\begin{remark}
The terminology ``unscreenable'' reflects the physics: the boundary mode $v$ cannot be
cancelled by any linear combination of interior modes.  It is \emph{topologically protected}
by the signed-incidence structure.  In the operator-factorization language of
\cite{Fontes2026I,Fontes2026II}, this is the analogue of the projection $\Pi$ isolating the
innovation: the component of the observable that survives conditioning on the interior.
\end{remark}

\subsection{Exact Schur complement for the quadratic core}

The Hessian of $V^{(0)}$ has block structure
\begin{equation}\label{eq:hessian}
\nabla^2 V^{(0)} \;=\; \beta c_2
\begin{pmatrix}
v^\top v & v^\top A \\
A^\top v & A^\top A
\end{pmatrix}.
\end{equation}

\begin{proposition}[Exact core Schur complement]\label{prop:core-schur}
Assume $A^\top A$ is invertible.  Then the Schur complement of $V_{yy}^{(0)}$ in $\nabla^2 V^{(0)}$ is
\begin{equation}\label{eq:core-schur}
\mathrm{Schur}(V^{(0)}) \;=\; \beta c_2\,v^\top(I-P_A)\,v.
\end{equation}
Under $A^\top v = 0$ (Lemma~\ref{lem:Atv0}) with $v=(1,\ldots,1)^\top\in\R^{n_p}$,
\begin{equation}\label{eq:schur-value}
\mathrm{Schur}(V^{(0)}) \;=\; n_p\,\beta\, c_2.
\end{equation}
In particular, for a $2\times 2$ tile with $n_p=4$ plaquettes, $\mathrm{Schur}(V^{(0)})=4\beta c_2$.
\end{proposition}

\begin{proof}
The cross-term vanishes: $V_{xy}^{(0)}=\beta c_2\,v^\top A = \beta c_2\,(A^\top v)^\top = 0$.
Therefore the Schur complement equals $V_{xx}^{(0)} - 0 = \beta c_2\,v^\top v = \beta c_2\,\norm{v}^2$.
\end{proof}

\begin{remark}
The exact vanishing of the cross-term is the operator-factorization at work:
integrating out interior variables does \emph{not} screen the boundary curvature.
This is the lattice analogue of the projection $\Pi$ in $(\mathrm{Id}-\E)=\delta\Pi D$
preserving the full innovation.
\end{remark}

\subsection{Robust Schur complement under perturbations}

The full potential $V = V^{(0)}+\Delta$ has a Hessian that deviates from the block-diagonal
structure.  We need a quantitative stability result.

\begin{assumption}[Interior coercivity]\label{ass:coerc}
There exists $\lambda_i>0$ such that $A^\top A\succeq \lambda_i I$.
\end{assumption}

\begin{assumption}[$C^2$ perturbation control]\label{ass:C2}
On the corridor $\norm{(x,y)}\le\delta$, the perturbation $\Delta$ satisfies
\[
\norm{\Delta_{xx}}\le K_{xx},\quad \norm{\Delta_{xy}}\le K_{xy},\quad \norm{\Delta_{yy}}\le K_{yy},
\]
in operator norm.
\end{assumption}

\begin{lemma}[Schur-complement stability]\label{lem:schur-stab}
Let $H, H_0$ be symmetric block matrices with $H_0$ having invertible lower-right block $C_0\succeq m_0 I$.
If $\norm{H-H_0}_{\mathrm{block}}\le(\eta_{xx},\eta_{xy},\eta_{yy})$ with $\eta_{yy}\le\frac{1}{2}m_0$,
and if $b_0$ denotes the off-diagonal block of $H_0$, then
\begin{equation}\label{eq:schur-stab}
\mathrm{Schur}(H) \;\ge\; \mathrm{Schur}(H_0)
- \eta_{xx} - \frac{4\eta_{xy}^2}{m_0} - \frac{2\norm{b_0}^2}{m_0^2}\,\eta_{yy}.
\end{equation}
\end{lemma}

\begin{proof}
Invertibility of the perturbed block: $\lambda_{\min}(C)\ge m_0-\eta_{yy}\ge\frac{1}{2}m_0$
by Weyl's inequality.  Expand $\mathrm{Schur}(H)-\mathrm{Schur}(H_0)=(a-a_0)-(b^\top C^{-1}b - b_0^\top C_0^{-1}b_0)$.
Split the quadratic-form difference:
\begin{align*}
b^\top C^{-1}b - b_0^\top C_0^{-1}b_0
&= (b-b_0)^\top C^{-1}(b-b_0) + 2(b-b_0)^\top C^{-1}b_0 + b_0^\top(C^{-1}-C_0^{-1})b_0.
\end{align*}
Bound each term using $\norm{C^{-1}}\le 2/m_0$ and the resolvent identity
$C^{-1}-C_0^{-1}=C^{-1}(C_0-C)C_0^{-1}$, then collect.
\end{proof}

\begin{theorem}[Robust boundary curvature]\label{thm:robust-schur}
Under Assumptions~\ref{ass:coerc}--\ref{ass:C2}, if $K_{yy}\le\frac{1}{2}\beta c_2\lambda_i$,
then
\begin{equation}\label{eq:robust-schur}
\mathrm{Schur}(V) \;\ge\; n_p\,\beta c_2 - K_{xx} - \frac{4K_{xy}^2}{\beta c_2\lambda_i}.
\end{equation}
In particular, if $K_{xx}=o(\beta)$ and $K_{xy}=o(\beta)$ as $\beta\to\infty$
(corridor scaling), then $\mathrm{Schur}(V)\ge\frac{n_p}{2}\,\beta c_2$ for $\beta$
sufficiently large.
\end{theorem}

\begin{proof}
Apply Lemma~\ref{lem:schur-stab} with $H=\nabla^2 V$, $H_0=\nabla^2 V^{(0)}$,
$m_0=\beta c_2\lambda_i$, and $b_0=0$ (by the topological cancellation).
The $b_0$-dependent term vanishes, giving~\eqref{eq:robust-schur}.
\end{proof}

\subsection{Strict $H^1$ contraction for the block transfer operator}

\begin{definition}[Block transfer operator]\label{def:T}
Let $\rho(x,y)=\exp(-V(x,y))\,\1_{\norm{(x,y)}\le\delta}$ and
$\nu_x(dy)\propto\rho(x,y)\,dy$ the conditional interior measure.  Define
\[
(Tf)(x) \;:=\; \E_{\nu_x}[f(X')],
\]
where $X'$ is the updated boundary variable after integrating out the interior.
\end{definition}

\begin{definition}[Reference operator]
Let $t_{\mathrm{eff}}:=(n_p\,\beta c_2)^{-1}$, motivated by the Schur-complement curvature.
Define $T_0:=P_{t_{\mathrm{eff}}}$, the heat semigroup on the boundary manifold at time
$t_{\mathrm{eff}}$.  If $\lambda_1>0$ is the spectral gap of $-\Delta$ on the boundary
manifold, then
\begin{equation}\label{eq:q0}
\norm{\nabla T_0 f}_{L^2} \;\le\; q_0\,\norm{\nabla f}_{L^2},\qquad
q_0 := e^{-\lambda_1 t_{\mathrm{eff}}} < 1.
\end{equation}
\end{definition}

\begin{assumption}[$H^1$ operator perturbation]\label{ass:HT}
$\norm{\nabla((T-T_0)f)}_{L^2}\le\varepsilon_T\,\norm{\nabla f}_{L^2}$ for some $\varepsilon_T\ge 0$.
\end{assumption}

\begin{assumption}[Interior Poincar\'e inequality]\label{ass:poinc}
$\Var_{\nu_x}(g)\le\kappa_{\mathrm{int}}^{-1}\,\E_{\nu_x}[\norm{\nabla_y g}^2]$ for all
smooth $g$, with $\kappa_{\mathrm{int}}\gtrsim\beta$ in the uniformly convex corridor.
\end{assumption}

\begin{lemma}[Covariance identity]\label{lem:cov-id}
For smooth $F(x,y)$ and any boundary derivative $D$,
\begin{equation}\label{eq:cov-id}
D\,\E_{\nu_x}[F] \;=\; \E_{\nu_x}[DF] - \Cov_{\nu_x}(F, DV).
\end{equation}
\end{lemma}

\begin{proof}
Differentiate $\E_{\nu_x}[F]=Z(x)^{-1}\int F\,e^{-V}\,dy$ by Leibniz' rule.  The
$DZ/Z^2$ term combines with the $-\int F(DV)e^{-V}/Z$ term to produce the covariance.
\end{proof}

\begin{lemma}[Covariance error bound]\label{lem:cov-bound}
Under the interior Poincar\'e inequality, the covariance correction contributes at most
$\eta\,\norm{\nabla f}_{L^2}$ with $\eta=C_{DV}/\kappa_{\mathrm{int}}^{1/2}$, where
$C_{DV}^2:=\sup_x\E_{\nu_x}[\norm{\nabla_y(DV)}^2]$.
\end{lemma}

\begin{theorem}[Strict block contraction]\label{thm:contraction}
Under the hypotheses of Theorem~\ref{thm:robust-schur} and
Assumptions~\ref{ass:HT}--\ref{ass:poinc},
\begin{equation}\label{eq:contraction}
\norm{\nabla(Tf)}_{L^2} \;\le\; q\,\norm{\nabla f}_{L^2},\qquad
q := q_0 + \varepsilon_T + \eta.
\end{equation}
If $q<1$, the block transfer operator is a strict $H^1$ contraction.
\end{theorem}

\begin{proof}
Decompose $\nabla(Tf)=\nabla(T_0 f)+\nabla((T-T_0)f)$, bound each by~\eqref{eq:q0}
and Assumption~\ref{ass:HT}, and add the covariance correction from
Lemma~\ref{lem:cov-bound}.  Triangle inequality gives~\eqref{eq:contraction}.
\end{proof}

\begin{corollary}[Exponential mixing under iteration]\label{cor:mixing}
If $q<1$, then $\norm{\nabla(T^d f)}_{L^2}\le q^d\,\norm{\nabla f}_{L^2}
=e^{-\mu d}\,\norm{\nabla f}_{L^2}$ with $\mu:=-\log q>0$.
\end{corollary}

%======================================================================
\section{From Local Contraction to Exponential Clustering}\label{sec:clustering}
%======================================================================

We now propagate the single-block estimate of Section~\ref{sec:schur} to global
correlation decay.

\subsection{Block lattice and iterated integration}

Let $\mathbb{L}$ denote the coarse lattice of blocks.  For a connected region
$\Omega\subset\mathbb{L}$, the iterated transfer operator integrates out blocks sequentially:
\[
T_\Omega := T_{B_N}\circ\cdots\circ T_{B_1}.
\]

\begin{definition}[Boundary Lipschitz seminorm]
$\Lip_\partial(G):=\bigl(\sum_{e\in\partial B}\sup_x\norm{\nabla_{x_e}G(x)}^2\bigr)^{1/2}$.
\end{definition}

\begin{lemma}[One-step influence contraction]\label{lem:onestep}
Under the block contraction (Theorem~\ref{thm:contraction} with $q<1$), there exists
$C_0<\infty$ such that $\Lip_\partial(T_B f)\le C_0\,q\,\Lip_\partial(f)$.
\end{lemma}

\begin{proof}
The $L^2$ contraction implies the Lipschitz estimate because the boundary marginal has a
smooth, strictly positive density on the corridor with uniformly bounded log-derivatives.
\end{proof}

\subsection{Exponential decay of boundary influence}

\begin{proposition}[Influence decay]\label{prop:influence}
Let $F$ be a bounded observable supported in blocks $\Delta\subset\Omega$.
For any two boundary conditions $x,x'$ on $\partial\Omega$,
\begin{equation}\label{eq:influence}
\big|\E_{\mu_\Omega(\cdot\mid x)}[F]-\E_{\mu_\Omega(\cdot\mid x')}[F]\big|
\;\le\; C_1\,\norm{F}_\infty\,e^{-\kappa\,\dist_{\mathbb{L}}(\Delta,\partial\Omega)}\,
d_\partial(x,x'),
\end{equation}
with $\kappa=-\log(C_0 q)>0$.
\end{proposition}

\begin{proof}
Integrate blocks layer-by-layer from the boundary inward.  After $n$ layers, the
Lipschitz seminorm of the effective observable has been reduced by $(C_0 q)^n$.
The mean-value theorem on the compact boundary converts Lipschitz bounds to
pointwise differences.
\end{proof}

\subsection{Exponential clustering of truncated correlations}

\begin{theorem}[Clustering from block contraction]\label{thm:clustering}
Let $\mu_\beta$ be the infinite-volume Gibbs measure in the corridor regime
(existence follows from Proposition~\ref{prop:influence} via Dobrushin uniqueness).
For bounded local observables $F,G$ supported in block sets $\Delta_F,\Delta_G$,
\begin{equation}\label{eq:cluster}
\big|\Cov_{\mu_\beta}(F,G)\big|
\;\le\; C_2\,\norm{F}_\infty\,\norm{G}_\infty\,
e^{-m\,\dist_{\mathbb{L}}(\Delta_F,\Delta_G)},
\end{equation}
where $m>0$ depends only on the corridor parameters and $q$.
\end{theorem}

\begin{proof}
Choose a buffer annulus $\Omega$ separating the supports.  Condition on the exterior
$\cF_{\Omega^c}$; by the DLR property, this is equivalent to conditioning on random
boundary data on $\partial\Omega$.  Apply Proposition~\ref{prop:influence} to bound
the fluctuation of $\E[F\mid\cF_{\Omega^c}]$ across the buffer, then integrate against $G$.
\end{proof}

%======================================================================
\section{Spectral Gap from Clustering and Reflection Positivity}\label{sec:gap}
%======================================================================

\subsection{Reflection positivity and transfer matrices}

\begin{assumption}[Reflection positivity]\label{ass:RP}
The lattice measure $\mu_\beta$ is reflection positive with respect to a coordinate
hyperplane, yielding a Hilbert space $\cH$, a self-adjoint transfer matrix
$\mathsf{T}=e^{-a\mathsf{H}}$ with $\mathsf{H}\ge 0$, and a vacuum vector $\Omega$.
\end{assumption}

\begin{theorem}[Spectral gap from exponential clustering]\label{thm:gap}
Under Assumption~\ref{ass:RP}, if truncated correlations satisfy~\eqref{eq:cluster} with
rate $m_{\mathrm{phys}}>0$, then the Hamiltonian $\mathsf{H}$ has a spectral gap at least
$m_{\mathrm{phys}}$ above the vacuum.
\end{theorem}

\begin{proof}
For a mean-zero time-local observable $F$, reflection positivity gives
$\ip{F}{e^{-t\mathsf{H}}F}_\cH=\E[F\,\Theta(\tau_t F)]$.
By the spectral theorem, this equals $\int_0^\infty e^{-t\lambda}\,d\mu_F(\lambda)$.
Exponential clustering forces this Laplace transform to decay as $e^{-m_{\mathrm{phys}}\,t}$,
which is only possible if $\mu_F$ is supported on
$\{0\}\cup[m_{\mathrm{phys}},\infty)$.  Since $F$ is arbitrary and mean-zero,
$\sigma(\mathsf{H})\cap(0,m_{\mathrm{phys}})=\emptyset$.
\end{proof}

\begin{corollary}[Mass gap from block contraction]\label{cor:gap}
Under Assumptions~\ref{ass:RP} and the block contraction (Theorem~\ref{thm:contraction}
with $q<1$), there exists $m_{\mathrm{phys}}>0$ such that the lattice theory has a spectral
gap $\ge m_{\mathrm{phys}}$.
\end{corollary}

\subsection{Persistence in continuum limits}

\begin{assumption}[Uniform contraction along a scaling family]\label{ass:uniform}
Along a trajectory $\beta=\beta(a)\to\infty$ as $a\downarrow 0$, the contraction constant
$q$ in Theorem~\ref{thm:contraction} satisfies $q<1$ uniformly in $a$.
\end{assumption}

\begin{theorem}[Continuum spectral gap]\label{thm:continuum-gap}
Under Assumptions~\ref{ass:RP}, \ref{ass:uniform}, and existence of subsequential continuum
limits of Schwinger functions, the limiting Euclidean theory satisfies exponential clustering
with rate $m_\star>0$.  The Osterwalder--Schrader reconstructed Hilbert space and Hamiltonian
have a spectral gap at least $m_\star$.
\end{theorem}

\begin{proof}
Uniform $q<1$ gives a uniform clustering rate $m_{\mathrm{phys}}$ independent of $a$
(Theorem~\ref{thm:clustering}).  Subsequential limits of bounded Schwinger functions preserve
the exponential decay bound by dominated convergence.  The OS reconstruction theorem
(cf.~\cite{OsterwalderSchrader1,OsterwalderSchrader2,GlimmJaffe}) produces a Hilbert space
with Hamiltonian $H\ge 0$, and the argument of Theorem~\ref{thm:gap} applies to the
continuum Schwinger functions.
\end{proof}

%======================================================================
\section{Uniqueness from Contraction: the Abstract $Ax=0$ Argument}\label{sec:uniqueness}
%======================================================================

The contraction mechanism yields not only decay estimates but a powerful \emph{uniqueness}
principle.  This section formulates the abstract version.

\subsection{Functional equations with contractive propagators}

Suppose the observables of the lattice theory satisfy a system of functional equations
(e.g.\ Schwinger--Dyson identities, loop equations, or lattice Ward identities) of the form
\begin{equation}\label{eq:func-eq}
\Phi(\{W_\alpha\}_{\alpha\in\mathcal{I}}) \;=\; 0,
\end{equation}
where $\{W_\alpha\}$ is a family of expectation values and $\Phi$ is a (possibly nonlinear)
operator.

\begin{definition}[Contractive linearization]
We say the system~\eqref{eq:func-eq} has a \emph{contractive linearization at rate $m$}
if the linearized operator $L:=D\Phi$ at any solution satisfies: for any two solutions
$W,W'$ of~\eqref{eq:func-eq},
\begin{equation}\label{eq:contractive}
\norm{W_\alpha - W'_\alpha}
\;\le\; C\sum_{\beta:\,\beta\ne\alpha}
e^{-m\,d(\alpha,\beta)}\,\norm{W_\beta - W'_\beta}
\end{equation}
in a suitable norm, where $d(\alpha,\beta)$ measures the separation of the observables.
\end{definition}

\begin{theorem}[Uniqueness from contraction plus boundary data]\label{thm:uniqueness}
Suppose the system~\eqref{eq:func-eq} has a contractive linearization at rate $m>0$,
and suppose additionally that boundary data (initial conditions, normalization, or
small-parameter asymptotics) fix the value of $W_\alpha$ for all $\alpha$ in a
``generating'' subset.  Then the system has at most one solution.
\end{theorem}

\begin{proof}
Let $W,W'$ be two solutions with the same boundary data, and let $\delta_\alpha:=W_\alpha-W'_\alpha$.
Then $\delta$ satisfies the homogeneous system
\[
\delta_\alpha \;=\; \sum_\beta L_{\alpha\beta}\,\delta_\beta,
\]
with $\delta_\alpha=0$ for generating $\alpha$ (by the shared boundary data).  The contractive
bound~\eqref{eq:contractive} with rate $m>0$ implies that the operator $L$ has spectral
radius strictly less than~1 in the weighted norm $\norm{\delta}_m:=\sup_\alpha e^{m\,|\alpha|}\,|\delta_\alpha|$.
Hence the only solution to $\delta = L\delta$ is $\delta\equiv 0$.
\end{proof}

\begin{remark}
This is the abstract content of the ``$Ax=0\Rightarrow x=0$'' argument:
\emph{exponential clustering makes the linearized propagator contractive},
so the homogeneous equation has only the trivial solution,
and \emph{boundary data pins the solution branch}.
In the stochastic calculus companion papers, the analogous statement is that
contraction of the conditional expectation yields uniqueness of the
martingale representation.
\end{remark}

\subsection{Application: lattice Schwinger--Dyson hierarchies}

In the lattice setting, integration-by-parts (Haar measure invariance for gauge theories,
or the standard Gibbs/DLR identity for spin systems) produces exact Schwinger--Dyson
equations for expectations of products of local observables.  Under the exponential clustering
of Theorem~\ref{thm:clustering}, these equations have contractive structure: each
observable's expectation is determined by nearby observables with exponentially decaying
coupling.  A normalization condition (e.g.\ fixing the two-point function of a chosen
observable at a reference separation) serves as the boundary data.

\begin{corollary}[Uniqueness of the gauge-invariant sector]\label{cor:uniqueness}
Under the hypotheses of Theorem~\ref{thm:clustering} (exponential clustering with rate $m>0$)
and the assumption that lattice Schwinger--Dyson identities pass to the continuum, the
expectations of all gauge-invariant observables in the continuum limit are uniquely determined
by:
\begin{enumerate}[label=(\alph*)]
\item the mass gap $m_\star$,
\item the continuum functional equations,
\item a single normalization datum (e.g.\ a two-point function value or small-parameter
asymptotic).
\end{enumerate}
\end{corollary}

%======================================================================
\section{Application to Lattice Gauge Theory with Compact Gauge Group}\label{sec:application}
%======================================================================

We briefly indicate how the abstract framework specializes to four-dimensional lattice gauge
theory with gauge group $G=\mathrm{SU}(2)$ (or any compact Lie group).

\subsection{Specialization of the block structure}

The lattice is $a\Z^4$ with gauge field $U_e\in G$ on each oriented edge.  Blocks are
$2\times 2\times 2\times 2$ hypercubes.  After spanning-tree gauge fixing within each block,
the interior variables are the link variables on non-tree edges, and the boundary variables
are the links crossing block boundaries.

The constraint vector $\omega$ consists of plaquette holonomies.  In exponential coordinates
$U_e=\exp(X_e)$ near identity, the linearized plaquette map has the form~\eqref{eq:omega}
with $A$ the signed incidence matrix and $v=(1,1,1,1)^\top$ for a $2\times 2$ tile.
Lemma~\ref{lem:Atv0} applies: each interior link appears in exactly two plaquettes with
opposite orientation.

The quadratic core coefficient is $c_2$, the quadratic coefficient in the even expansion
of the Wilson plaquette potential $\frac{\beta}{2}(2-\mathrm{Re}\,\tr(U_p))$
$=\beta c_2\norm{X_p}^2+O(\norm{X_p}^4)$.  The evenness of the SU(2) expansion (no cubic
term) is crucial for the corridor control: it ensures the perturbation $\Delta$ satisfies
$K_{xx},K_{xy}=O(\beta\delta^2)+O(1)$, which is $o(\beta)$ when $\delta\to 0$.

The boundary manifold is a product of copies of $G$, and the spectral gap $\lambda_1$ of
$-\Delta$ on $G$ is known explicitly (e.g.\ $\lambda_1=2$ for SU(2) with the round metric
of diameter $\pi$).

\subsection{Verification of hypotheses}

\begin{enumerate}[label=(\roman*)]
\item \emph{Topological cancellation}: Lemma~\ref{lem:Atv0} holds by the combinatorics
of the $2\times 2$ tile.
\item \emph{Interior coercivity}: $A^\top A\succeq\lambda_i I$ with $\lambda_i>0$
after spanning-tree gauge fixing (the remaining plaquette constraints are linearly
independent on the interior links).
\item \emph{$C^2$ perturbation control}: Follows from BCH evenness and uniform corridor
bounds on exponential coordinates.
\item \emph{Interior Poincar\'e}: The corridor potential is uniformly convex with
constant $\gtrsim\beta$, so Brascamp--Lieb applies.
\item \emph{Reflection positivity}: Standard for the Wilson action on $a\Z^4$ with
appropriate boundary conditions \cite{OsterwalderSeiler}.
\end{enumerate}

The contraction constant $q=q_0+\varepsilon_T+\eta$ satisfies $q<1$ for $\beta$
sufficiently large and $\delta$ sufficiently small.  The resulting spectral gap is
$m_{\mathrm{phys}}>0$.

\subsection{Continuum limit}

Under the assumption that a scaling trajectory $\beta(a)\to\infty$ exists with uniform
corridor control (Assumption~\ref{ass:uniform}), Theorem~\ref{thm:continuum-gap} yields
a continuum Euclidean theory with a positive spectral gap.

The identification of this theory as a specific gauge theory (rather than an arbitrary
gapped theory) requires additional input: passage of lattice Ward identities to the
continuum and a nontriviality condition.  These are discussed in the companion notes
\cite{Fontes2026Ward,Fontes2026Loop} and are outside the scope of the present
operator-factorization framework.

%======================================================================
\section{Discussion}\label{sec:discussion}
%======================================================================

\subsection{The operator-factorization thread}

The central message of this paper is that the same ``factor, project, contract'' philosophy
that underlies Clark--Ocone representations and Bakry--\'Emery curvature in the stochastic
calculus setting also drives mass-gap extraction in lattice statistical mechanics.  The
specific instantiations differ:
\begin{center}
\begin{tabular}{lcc}
& \emph{Stochastic calculus} & \emph{Lattice blocks}\\[3pt]
\hline\\[-6pt]
Operator identity & $(\mathrm{Id}-\E)=\delta\Pi D$ & $\mathrm{Schur}(V)=V_{xx}-V_{xy}V_{yy}^{-1}V_{yx}$\\[2pt]
Topological input & Adaptedness / filtration & $A^\top v=0$\\[2pt]
Contraction & $\norm{(\mathrm{Id}-\E)f}\le q\norm{f}$ & $\norm{\nabla(Tf)}\le q\norm{\nabla f}$\\[2pt]
Global conclusion & Representation theorem & Spectral gap / mass gap\\[2pt]
Uniqueness & Martingale uniqueness & $Ax=0\Rightarrow x=0$
\end{tabular}
\end{center}
We view these as manifestations of a single structural principle rather than independent results.

\subsection{What this framework does and does not accomplish}

The framework provides a clean, deterministic pipeline:
\begin{quote}
\emph{Local contraction} $\Longrightarrow$ \emph{Exponential clustering}
$\Longrightarrow$ \emph{Spectral gap} $\Longrightarrow$ \emph{Uniqueness}.
\end{quote}
What it does \emph{not} supply is the local contraction estimate itself for any specific
model.  This requires model-specific analysis (corridor bounds, BCH expansions,
character-algebraic estimates) that is the subject of the companion papers.  Nor does the
framework address identification of the continuum limit with a specific field theory:
that requires Ward identities and nontriviality, which are logically independent of the
contraction mechanism.

The separation is deliberate: by isolating the abstract pipeline, we clarify that
the spectral gap is a \emph{deterministic consequence} of a single local estimate,
and the only genuinely difficult model-dependent step is establishing that estimate.

\subsection{Open directions}

\begin{enumerate}[label=(\arabic*)]
\item \emph{Explicit corridor verification}: Close the contraction estimate for SU(2)
lattice gauge theory with explicit constants, establishing $q<1$ for a concrete range
of $\beta$.
\item \emph{Higher-rank groups}: Extend the topological cancellation and corridor analysis
to SU($N$) for $N\ge 3$.
\item \emph{Non-gauge lattice systems}: The Schur-complement mechanism applies whenever
a block decomposition has the $A^\top v=0$ structure; explore applications to lattice
$\sigma$-models and other systems with topologically protected modes.
\item \emph{Quantitative continuum bounds}: Combine the contraction pipeline with quantitative
OS reconstruction to produce explicit bounds on the continuum spectral gap as a function
of the lattice contraction constant.
\end{enumerate}

%======================================================================
% References
%======================================================================

\bibliographystyle{alpha}
\begin{thebibliography}{99}

\bibitem[Bal84]{Balaban1}
T.~Balaban,
\newblock Propagators and renormalization transformations for lattice gauge fields~I.
\newblock \emph{Comm.\ Math.\ Phys.}, 95:17--40, 1984.

\bibitem[Bal85a]{Balaban2}
T.~Balaban,
\newblock Averaging operations for lattice gauge theories.
\newblock \emph{Comm.\ Math.\ Phys.}, 98:17--51, 1985.

\bibitem[Bal85b]{Balaban3}
T.~Balaban,
\newblock Spaces of regular gauge field configurations on a lattice and gauge fixing conditions.
\newblock \emph{Comm.\ Math.\ Phys.}, 99:75--102, 1985.

\bibitem[DS85]{DobrushinShlosman}
R.L.~Dobrushin and S.B.~Shlosman,
\newblock Constructive criterion for the uniqueness of Gibbs field.
\newblock In \emph{Statistical Physics and Dynamical Systems} (K\"osz\"eg, 1984),
Birkh\"auser, 1985.

\bibitem[FI]{Fontes2026I}
R.~Fontes,
\newblock Operator factorization and Clark--Ocone formulas for stable L\'evy processes.
\newblock Quijotic Research, 2025--2026.

\bibitem[FII]{Fontes2026II}
R.~Fontes,
\newblock Operator factorization and It\^o formulas for fractional Brownian motion.
\newblock Quijotic Research, 2025--2026.

\bibitem[FIII]{Fontes2026III}
R.~Fontes,
\newblock Bakry--\'Emery curvature via operator factorization.
\newblock Quijotic Research, 2025--2026.

\bibitem[FIV]{Fontes2026IV}
R.~Fontes,
\newblock Martingale innovations and non-semimartingale It\^o formulas.
\newblock Quijotic Research, 2025--2026.

\bibitem[FW]{Fontes2026Ward}
R.~Fontes,
\newblock Lattice Ward identities and their continuum passage.
\newblock Quijotic Research, 2026 (in preparation).

\bibitem[FL]{Fontes2026Loop}
R.~Fontes,
\newblock Continuum loop equations from lattice Schwinger--Dyson identities.
\newblock Quijotic Research, 2026 (in preparation).

\bibitem[GJ]{GlimmJaffe}
J.~Glimm and A.~Jaffe,
\newblock \emph{Quantum Physics: A Functional Integral Point of View.}
\newblock Springer, 2nd edition, 1987.

\bibitem[OS73]{OsterwalderSchrader1}
K.~Osterwalder and R.~Schrader,
\newblock Axioms for {E}uclidean {G}reen's functions.
\newblock \emph{Comm.\ Math.\ Phys.}, 31:83--112, 1973.

\bibitem[OS75]{OsterwalderSchrader2}
K.~Osterwalder and R.~Schrader,
\newblock Axioms for {E}uclidean {G}reen's functions~{II}.
\newblock \emph{Comm.\ Math.\ Phys.}, 42:281--305, 1975.

\bibitem[OSe78]{OsterwalderSeiler}
K.~Osterwalder and E.~Seiler,
\newblock Gauge field theories on a lattice.
\newblock \emph{Ann.\ Physics}, 110:440--471, 1978.

\end{thebibliography}

\end{document}
